\documentclass{article}
\usepackage{graphicx}
\usepackage{todonotes}

\usepackage[T1]{fontenc}
\usepackage{libertinus,libertinust1math}

\usepackage[a4paper, total={6.5in, 9.5in}]{geometry}

\usepackage[style=alphabetic]{biblatex}
\addbibresource{references.bib}

\title{parallel. Forced and time-dependent}
\author{David Martin de Diego}
\date{December 2025}


\begin{document}

\maketitle

\section{Dissipative forces}

The method in the \cite{ferraroParallelIterativeMethod2025} works for equations of the style 

\[ \Diff_2 L_d (q_{k-1}, q_k) + \Diff_1 L_d (q_k, q_{k+1}) = 0 \]

But what if we have a force?

\[ \Diff_2 L_d (q_{k-1}, q_k) + \Diff_1 L_d (q_k, q_{k+1}) = F^+(q_{k-1}, q_k) + F^-(q_k, q_{k+1}) \]

This is a dissipative force \todo{creo, vamos}. 

We want to figure out when solving problems with dissipative forces using this method \todo{más o menos} converges. At a glance, we expect that it should always converge, because you can just modify the lagrangian by adding a $e^{-Ft}$ factor to convert it to a non-dissipative problem \todo{no se si esto siempre pasa}. 

However, how can we properly verify if it converges?

\subsection{Duplicating dimensions}

Given some state\todo{es state space?} space of dimension $n$, we can pose a new problem of dimension $2n$ where we have $q_k$ for $k = 1,...,n$ but also $Q_k$ for $k = 1,...,n$ and with no force. When $q_k = Q_k$ for all $k$, then $q_k$ is a solution to the original (forced) problem.

And here we can use the converge criteria we have from \cite{ferraroParallelIterativeMethod2025}! Namely, we need the following matrix to be positive-definite:

\[
    H(\mathring{q}_d^{[\gamma-1]}) = \begin{pmatrix}
        \mathcal{D}_1 & \mathcal{C}_1 \\
        \mathcal{C}_1^T & \mathcal{D}_2 & \mathcal{C}_2 \\
        & \ddots & \ddots & \ddots & \\
        & & \mathcal{C}_{N-3}^T & \mathcal{D}_{N-2} & \mathcal{C}_{N-2} \\ 
        & & & \mathcal{C}_{N-2}^T & \mathcal{D}_{N-1} \\ 
        
    \end{pmatrix}
\]

where 

\begin{align*}    
    \mathcal{D}_k &= \Diff_{22} L_d (q_{k-1}^{[\gamma - 1]}, q_k^{[\gamma - 1]}) + \Diff_{1 1} L_d (q_k^{[\gamma - 1]}, q_{k+1}^{[\gamma - 1]}) \\
    \mathcal{C}_k &= \Diff_{12} L_d (q_k^{[\gamma - 1]}, q_{k+1}^{[\gamma - 1]})
\end{align*}

Note that this condition is not necessary, but it is sufficient \todo{is it?}.

The superscripts indicate iteration index, but we're not interested in that, so we can just straight up ignore them. What we want to see is that if the non-forced version of the equation converges, then if this will converge too. 

Let's start with a simple example, with a $1$-dimensional state space. The Lagrangian we are going to take is the following (note that $D$ here is the dissipation force, not $\Diff$, the derivative):

\todo{habría que cambiar $q_0 \to q_k, q_1 \to q_{k+1}$?}

\[ L_d (q_0, Q_0, q_1, Q_1) = - \frac{D h \left(Q_{0} - q_{0}\right) \left(\frac{- Q_{0} + Q_{1}}{h} + \frac{- q_{0} + q_{1}}{h}\right)}{2} + \frac{\left(- Q_{0} + Q_{1}\right)^{2}}{2 h} - \frac{\left(- q_{0} + q_{1}\right)^{2}}{2 h} \]

\[ L_d (q_0, Q_0, q_1, Q_1) =
            \frac{h}{2} \left(\frac{Q_1 - Q_0}{h}\right)^2
            - \frac{h}{2} \left(\frac{q_1 - q_0}{h}\right)^2
            - \frac{h}{2} \left(D \left(\frac{Q_1 - Q_0}{h} + \frac{q_1 - q_0}{h}\right) (Q_0 - q_0)\right)
\]

\subsubsection{Non-dissipative case}

We are going to first look at what happens, first, if we solve the non forced equation, pretending we only have the $n$-dimensional (in this case, 1-dimensional) state space. We do this by taking the derivatives with respect to $q$ and $Q$ independently, and by the nature of the construction we obtain the same result. 

\textcolor{red}{En verdad estoy confundido con qué estoy haciendo aquí exactamente. Es verdad que siempre va a salir lo mismo?}

The first derivative with respect to $q_0$ is

\[
\frac{1}{h} \left(
\frac{D h \left(Q_{1} - 2 q_{0} + q_{1}\right)}{2} - q_{0} + q_{1}
\right)
\]

The second derivative with respect to $q_0$ is

\[
    - D - \frac{1}{h}
\]

We get the same with opposite signs for $Q$. Notice that it doesn't depend on $q$ nor $Q$. Therefore, all $\mathcal{D}_k$s and $\mathcal{C}_k$s are going to be identical. If we compute $\Dif_{12}$ and $\Dif_{22}$ we are left with submatrices:

\begin{align}    
    \mathcal{D}_k & = - D - \frac{2}{h} \\
    \mathcal{C}_k & = \frac{D}{2} + \frac{1}{h}
\end{align}

So then, if we construct the matrix and we take a factor of $\frac{D}{2} + \frac{1}{h}$, we get


\[
    H(q) = \left(\frac{D}{2}+ \frac{1}{h}\right) \begin{pmatrix}
        -2 & 1 \\
        1 & -2 & 1 \\
        & \ddots & \ddots & \ddots & \\
        & & 1 & -2 & 1 \\ 
        & & & 1 & -2 \\  
    \end{pmatrix}
\]

The inside matrix is negative definite so the whole matrix is positive definite if the factor is negative, i.e., if

\begin{align}
    \frac{D}{2} + \frac{1}{h} & \le 0 \\
    % h & \le - \frac{2}{D} \\
    Dh & \le -2
\end{align}


This makes sense, in the sense that $D$ first needs to be negative to be an actual dissipation force (otherwise the energy (?) exponentially increases and it leads to instability). However, the timestep needs to be small if the dissipation is small, which does seem a bit weird to me.

\subsubsection{Dissipative case}

Now, if we compute the Hessian using ``extended'' state space, using a $L(q_0, q_1)$-type Lagrangian with inputs, $\begin{pmatrix} q_0 \\ Q_0 \end{pmatrix}$ and $\begin{pmatrix} q_1 \\ Q_1 \end{pmatrix}$. The Hessian is

\[
    \left[\begin{matrix}- D - \frac{1}{h} & 0 & \frac{D}{2} + \frac{1}{h} & \frac{D}{2}\\0 & D + \frac{1}{h} & - \frac{D}{2} & - \frac{D}{2} - \frac{1}{h}\\\frac{D}{2} + \frac{1}{h} & - \frac{D}{2} & - \frac{1}{h} & 0\\\frac{D}{2} & - \frac{D}{2} - \frac{1}{h} & 0 & \frac{1}{h}\end{matrix}\right]
\]

so,

\begin{align}    
    \mathcal{D}_k & = \left[\begin{matrix}- D - \frac{2}{h} & 0\\0 & D + \frac{2}{h}\end{matrix}\right] \\
    \mathcal{C}_k & = \left[\begin{matrix}\frac{D}{2} + \frac{1}{h} & \frac{D}{2}\\- \frac{D}{2} & - \frac{D}{2} - \frac{1}{h}\end{matrix}\right]
\end{align}

Again, this is independent of $k$. However, this time it's not so easy to write out the matrix because it doesn't factor out cleanly into one matrix, but we can write it as a sum. Namely,

\begin{align*}    
    H(q) & = \begin{pmatrix}
        \mathcal{D}_1 & \mathcal{C}_1 \\
        \mathcal{C}_1^T & \mathcal{D}_2 & \mathcal{C}_2 \\
        & \ddots & \ddots & \ddots & \\
        & & \mathcal{C}_{N-3}^T & \mathcal{D}_{N-2} & \mathcal{C}_{N-2} \\ 
        & & & \mathcal{C}_{N-2}^T & \mathcal{D}_{N-1} \\  
    \end{pmatrix} \\
    & = \begin{pmatrix}
        - D - \frac{2}{h} & 0 & \frac{D}{2} + \frac{1}{h} & \frac{D}{2} & 0 & 0\\
        0 & D + \frac{2}{h} & - \frac{D}{2} & - \frac{D}{2} - \frac{1}{h} & 0 & 0\\
        \frac{D}{2} + \frac{1}{h} & - \frac{D}{2} & - D - \frac{2}{h} & 0 & \frac{D}{2} + \frac{1}{h} & \frac{D}{2}\\
        \frac{D}{2} & - \frac{D}{2} - \frac{1}{h} & 0 & D + \frac{2}{h} & - \frac{D}{2} & - \frac{D}{2} - \frac{1}{h}\\
        & \ddots & \ddots & \ddots & \ddots \\
 0 & 0 & \frac{D}{2} + \frac{1}{h} & - \frac{D}{2} & - D - \frac{2}{h} & 0\\0 & 0 & \frac{D}{2} & - \frac{D}{2} - \frac{1}{h} & 0 & D + \frac{2}{h}
    \end{pmatrix} \\
    & = \begin{pmatrix}
        - D & 0 & \frac{D}{2} & \frac{D}{2} & 0 & 0\\
        0 & D & - \frac{D}{2} & - \frac{D}{2}  & 0 & 0\\
        \frac{D}{2} & - \frac{D}{2} & - D & 0 & \frac{D}{2}  & \frac{D}{2}\\
        \frac{D}{2} & - \frac{D}{2}  & 0 & D & - \frac{D}{2} & - \frac{D}{2} \\
        & \ddots & \ddots & \ddots & \ddots \\
 0 & 0 & \frac{D}{2} & - \frac{D}{2} & - D & 0\\0 & 0 & \frac{D}{2} & - \frac{D}{2}  & 0 & D
    \end{pmatrix} + \begin{pmatrix}
        - \frac{2}{h} & 0 & \frac{1}{h} & 0 & 0 & 0\\
        0 & \frac{2}{h} & 0 &- \frac{1}{h} & 0 & 0\\
         \frac{1}{h} & 0 & - \frac{2}{h} & 0 &  + \frac{1}{h} & 0\\
       0 &  - \frac{1}{h} & 0 & \frac{2}{h} & 0 & - \frac{1}{h}\\
        & \ddots & \ddots & \ddots & \ddots \\
 0 & 0 & \frac{1}{h} & 0 & - \frac{2}{h} & 0\\0 & 0 & 0 & - \frac{1}{h} & 0 & \frac{2}{h}
    \end{pmatrix} \\
    & = \frac{-D}{2} \begin{pmatrix}
        2 & 0 & -1 & -1 & 0 & 0\\
        0 & -2 & 1 & 1  & 0 & 0\\
        -1 & 1 & 2 & 0 & -1  & -1\\
        -1 & 1  & 0 & -2 & 1 & 1 \\
        & \ddots & \ddots & \ddots & \ddots \\
 0 & 0 & -1 & 1 & 2 & 0\\0 & 0 & -1 & 1  & 0 & -2
    \end{pmatrix} - \frac{1}{h} \begin{pmatrix}
        2  & 0 & -1 & 0 & 0 & 0\\
        0 & -2  & 0 &1 & 0 & 0\\
         -1 & 0 & 2  & 0 &  + -1 & 0\\
       0 &  1 & 0 & -2  & 0 & 1\\
        & \ddots & \ddots & \ddots & \ddots \\
 0 & 0 & -1 & 0 & 2 & 0\\0 & 0 & 0 & 1 & 0 & -2 
    \end{pmatrix} 
\end{align*}

It seems that this matrix, though, is never positive definite.

In fact, in the joint matrix, if we use Sylvester's criterion, the first minor would give the $-D - \frac{2}{h} > 0$ condition which is the one we had in the non-dissipative case ($Dh < -2$); but for the second already you get 

\[\begin{aligned}
    \left(-D - \frac{2}{h}\right) \left(D + \frac{2}{h}\right) & > 0 \\
    -\left(D + \frac{2}{h}\right) \left(D + \frac{2}{h}\right) & > 0 \\
    -\left(D + \frac{2}{h}\right)^2 &> 0 \\
\end{aligned}\]

which can never happen. Therefore, this matrix is \emph{not} p.d..

\subsubsection{Constrained positive definiteness}

A matrix is p.d. if, for all $x$, $x^T A x > 0$. However, here our $x$s are 

\[
    \begin{pmatrix}
        q_0 \\
        Q_0 \\
        q_1 \\
        Q_1 \\
        q_2 \\
        Q_2 \\
        \vdots
    \end{pmatrix}
\]

but perhaps we're giving too much freedom to this vector spanning the entirety of $\mathbb{R}^{2n}$, since we only care about solutions where $q_0 = Q_0$. So what if we allow $x$ to be only in $\mathbb{R}^{2n} \land x_{2i} = x_{2i+1}$ for $i = 0, ..., n-1$? 

Let's call these type of matrices \emph{constrained p.d.} or CPD. Is $H$ CPD?

Well, if we carry out this calculation for this matrix on the computer we get this monstrosity...

\[
\begin{aligned}    
x_{0} \left(- \frac{D x_{1}}{2} + x_{0} \left(- D - \frac{2}{h}\right) + x_{1} \left(\frac{D}{2} + \frac{1}{h}\right)\right) \\ - x_{0} \left(- \frac{D x_{1}}{2} - x_{0} \left(D + \frac{2}{h}\right) - x_{1} \left(- \frac{D}{2} - \frac{1}{h}\right)\right) \\ - x_{1} \left(\frac{D x_{0}}{2} - \frac{D x_{2}}{2} - x_{0} \left(- \frac{D}{2} - \frac{1}{h}\right) - x_{1} \left(D + \frac{2}{h}\right) - x_{2} \left(- \frac{D}{2} - \frac{1}{h}\right)\right) \\ + x_{1} \left(\frac{D x_{0}}{2} - \frac{D x_{2}}{2} + x_{0} \left(\frac{D}{2} + \frac{1}{h}\right) + x_{1} \left(- D - \frac{2}{h}\right) + x_{2} \left(\frac{D}{2} + \frac{1}{h}\right)\right) \\ - x_{2} \left(\frac{D x_{1}}{2} - \frac{D x_{3}}{2} - x_{1} \left(- \frac{D}{2} - \frac{1}{h}\right) - x_{2} \left(D + \frac{2}{h}\right) - x_{3} \left(- \frac{D}{2} - \frac{1}{h}\right)\right) \\ + x_{2} \left(\frac{D x_{1}}{2} - \frac{D x_{3}}{2} + x_{1} \left(\frac{D}{2} + \frac{1}{h}\right) + x_{2} \left(- D - \frac{2}{h}\right) + x_{3} \left(\frac{D}{2} + \frac{1}{h}\right)\right) \\ - x_{3} \left(\frac{D x_{2}}{2} - x_{2} \left(- \frac{D}{2} - \frac{1}{h}\right) - x_{3} \left(D + \frac{2}{h}\right)\right) \\ + x_{3} \left(
\frac{D x_{2}}{2} + x_{2} \left(\frac{D}{2} + \frac{1}{h}\right) + x_{3} \left(- D - \frac{2}{h}\right)\right)
\end{aligned}
\]

...which turns out to be exactly $0$. This might happen in general when you do this duplicating of dimensions to convert a dissipative problem into a non-dissipative one, because this seems to unlikely.

Still, not constrained positive definite. If anything, constrained semi-positive definite, but really we've just shown that the solution vectors are the null space of this matrix, for some reason. Maybe this is interesting.

\subsection{Time-dependent Lagrangian}

Theoretically, this is all very easy to extend to time-dependent Lagragians, but there are some complications with the indices.

First, note that $k$ indicates the time, since $q_k$ are points of the path in time. Then, the discrete Lagrangian now has an index $k$ like $L_k (q_k, q_{k+1})$. Note that here the $k$ of the $q$ \textbf{is} relevant.

The discrete Euler-Lagrange equations now become: 

\[ \Diff_2 L_{k-1} (q_{k-1}, q_k) + \Diff_1 L_k (q_k, q_{k+1}) = 0 \]

and with a force,

\[ \Diff_2 L_{k-1} (q_{k-1}, q_k) + \Diff_1 L_k (q_k, q_{k+1}) = F \]

Honestly, I'm not sure why I said it was confusing lol.

\subsection{Exponential decay}

Ok so the explicit discrete Lagrangian with a dissipative force $D$ is:

\[
    L_k(q_k, q_{k+1}) = h e^{D h k} \left(\frac{1}{2} \left(\frac{q_{k+1} - q_k}{h}\right)^2\right)
\]

Given the Hessian components:

\begin{align*}    
    \mathcal{D}_k &= \frac{2}{h} e^{D h k} \\
    \mathcal{C}_k &= - \frac{1}{h} e^{D h k}
\end{align*}

we can compute the big hessian:

\[
    H(q) = \frac{1}{h} \left[\begin{matrix}2 & -1 & 0 & 0 & 0\\-1 & 2 e^{D h} & - e^{D h} & 0 & 0\\0 & - e^{D h} & 2 e^{2 D h} & - e^{2 D h} & 0\\0 & 0 & - e^{2 D h} & 2 e^{3 D h} & - e^{3 D h}\\0 & 0 & 0 & - e^{3 D h} & 2 e^{4 D h}\end{matrix}\right]
\]

I don't immediately see if this matrix is positive definite.

\subsubsection{Sylvester's criterion}

But what if we look at Sylvester's criterion? Iff every leading principal minor is positive, then the matrix is p.d.. So, we need the following to be positive (skipping the division over power of $h$, since that preserves positivity):

\begin{align*}
    & 2 \\
    & 4 e^{D h} - 1 \\
    & \left(8 e^{D h} - 4\right) e^{2 D h} \\
    & \left(16 e^{2 D h} - 12 e^{D h} + 1\right) e^{4 D h} \\
    & \left(32 e^{2 D h} - 32 e^{D h} + 6\right) e^{8 D h} \\
\end{align*}

\begin{enumerate}
    \item $2 > 0$ imposes no condition
    \item $4 e^{D h} - 1 > 0$ implies $D h > \log \frac{1}{4} \approx -1.38$
    \item $\left(8 e^{D h} - 4\right) e^{2 D h} > 0$ implies $8 e^{D h} - 4 > 0 \Rightarrow D h > \log \frac{1}{2} \approx -0.693$
\end{enumerate}

From this we see that the converging criteria don't seem to be the same than for the non-dissipative case. There, we needed $Dh \leq -2$, but here we have the equality the wrong way around.

\subsection{Continuous Lagrangian}

So, we are essentially solving the system

\[
    \ddot{q} + D \dot{q} = 0
\]

which has Lagrangian

\[
    L = e^{D t} \left(\frac{1}{2} \dot{q}^2\right)
\]

what can we see from this?



\break
\printbibliography

\end{document}
